\chapter{Introduction}

\section{Why microscopy?}
Microscopy of biological specimens is used in biotechnology to
visualize small specimens (including their internal structures) that
cannot be seen with the naked eye. For this, some kind of interaction
between the specimens and a source of energy must exist (the specimen
must be irradiated), an action that usually degrades the specimen and
for this reason, the radiation is minimized (in radiation power or
``dwell'' time) generating low-SNR images. In this report, we describe,
evaluate, and compare several image denoising algorithms commonly used
in microscopy of biological specimens.

\section{Image acquisition in microscopy}
%{{{

There are several tecniques to capture microscopy images:

\begin{enumerate}
\item In \textbf{\gls{LM}}, light passes through the specimen, and this
  light is then magnified by an objective and ocular lenses to form an
  observable image. Fluorescence microscopy\footnote{The target of
  interest is ``tinted'' with fluorophores that upon ilumination with
  a specific wavelength of excitation light, the attached fluorophores
  emit photons.} and confocal microscopy are two \gls{LM} tecniques,
  where the specimen(s) emit light after have being excited by a
  source of light. Using \gls{LM} we can typical reach a resolution of
  up to 200 nm.

\item In \textbf{\gls{EM}}, we use a beam of electrons that can pass
  through the sample to discover the internal structure (\gls{TEM}) or
  bounce on the sample, in this case to analize the surface topology
  (\gls{SEM}) \cite{timischl2012statistical}. X-rays and ions can also
  be used (the smaller the wavelength of the radiation, the higher the
  resolution). In general, compared to \gls{LM}, \gls{EM} allows much
  higher resolution (down to sub-nanometer scale).

\item \textbf{\gls{SPM}} uses a nano-scale mechanical probe that maps
  the surface of a specimen. \gls{SPM} achieves sub-nanometer
  resolution, allowing for the imaging of surfaces at the atomic
  scale. The \gls{AFM} and the \gls{STM} are classified as \gls{SPM}
  techniques.

\end{enumerate}
  
%}}}

\section{Tomography}
%{{{

Tomography refers to the (usually computational) reconstruction of 3D
volumes from a series of 2D projections. Depending on the imaging
modalities we can distinguish bewtween:
\begin{enumerate}
\item \textbf{\gls{ET}}: Used mainly in cell biology and materials
  science. The samples are normally resin-embedded. Achieve resolutions
  btween 1 and 10 nm.
\item \textbf{\gls{cryo-ET}}: Used in structural biology (e.g.,
  visualizing viruses, ribosomes, and organelles). To minimize the
  damages in the samples they are frozen-hydrated using a
  flash-freezing technique, called \emph{vitrification}, which
  suspends the specimens in a thin layer of amorphous, non-crystaline
  ice\footnote{Which prevents the formation of damaging ice crystals
  and therefore preserving the sample's structure.}. \gls{cryo-ET} achieves
  resolutions range between 3 and 5 nm. X-rays can replace
  electron-beams.
\item \textbf{\gls{OPT}}: Normally based in light microscopy,
  \gls{OPT} is a non-destructive 3D imaging technique used to observe
  transparent samples like embryos and organoids. The typical
  resolutions range between 5 and 50 $\mu$m.
\item \textbf{Fluorescence Tomography}: Used to mapping 3D
  fluorescence distribution in tissues, oftenly for molecular imaging
  in live animals or cleared tissues. Resolution ranges between 10 and
  100 $\mu$m.
\item \textbf{\gls{SXT}}: 3D imaging of intact cells with
  higher penetration than electrons or light. Resolution between 30
  and 50 nm.
\item \textbf{\gls{PAT}}: Light excitation + ultrasound
  detection. Used in deep-tissue imaging of functional contrast (like
  blood oxygenation) in live animals. Resolution between 50 and 300
  $\mu$m.
\end{enumerate}

Therefore, a tomogram is a 3D digital signal (volume) generated by a
tomography technique.

%}}}

\section{Sources of distortion}

The typical distortions that can be fould in miscroscopy include:
\begin{enumnerate}
  \item The
    effects of camera and photon/electron noise,
  \item the blur of the optical
    point-spread-function.
  \item the resolution loss due to sampling, and
  \item the aberrations induced by the sample itself.
\end{enumerate}
Although we have a very detailed understanding of the nature and
physics behind contributing distortions, it is hard to compute the
inverse, i.e., generate a reconstructed image by only observing the
corrupted images acquired by microscopes.\footnote{For this reason,
these type of problems are classified as \emph{non-invertible}.} The
inverse of an observed image $\mathbf{X}$ is usually not uniquely
defined, meaning that multiple undistorted images could have given
rise to the corrupted version we observed \cite{buchholz2019content}.

\section{Sources of noise}

In microscopy imaging there are tho main sources of noise
\cite{meiniel2018denoising,zhou2020wirtinger}:
\begin{enumerate}
\item \textbf{Photon} (also called \textbf{shot}) \textbf{noise},
  which is generated by the own stochastic nature of the light which
  produces fluctuations in the number of phothons
  sensed. \footnote{The light travels in discrete packets of energy
  called photons, and these photons do not arrive at the sensor in a
  perfectly steady stream; they arrive randomly.} This type of noise
  is usually modeled as a Poisson distribution $\mathcal{P}(\lambda)$,
  where $\lambda$ represents the average dark flux
\item \textbf{Dark noise}, which corresponds to the noise generated by
  the thermal agitation of the electrons in the circuits of the camera
  (CCD, amplifiers, etc.). Thermal noise is modeled as zero-mean
  additive white Gaussian noise.
\end{enumerate}


